\section{평가 방법론 및 결과}
\label{sec:evaluation_methodology}

본 절은 시스템의 탐지 및 대응 성능을 어떻게 측정하였는지, 그리고 그 결과가 무엇을 의미하는지를 서술한다. 인용된 모든 수치는 \texttt{docs/evidence/} 하위의 실제 실행 산출물에 근거하며, 측정할 수 없는 항목은 미측정 또는 \texttt{null}로 표기한다. 평가는 다섯 개의 상호 보완적 실험으로 구성되며, 다음과 같다: (1) 결정론적 오라클 채점기 기반의 객관적 탐지 지표, (2) 도구 미사용 방어군 (blind group) 과 도구 결합형 지휘군 (agentic commander group) 간의 대조 실험, (3) 세 모델의 벤치마크 평가 (bake-off), (4) 3개 도메인 대상의 통합 시뮬레이션 조율 캠페인 (real-game), (5) D5 탐지기의 경험적 탐지 경계 분석.

\subsection{핵심 탐지 지표에서 LLM 자기평가 순환을 배제}
\label{subsec:eval_methodology}

% 초기 평가는 에피소드별로 LLM이 자신의 진단 결과를 $0\sim100$점으로 직접 채점하는 방식을 사용하였다. 그러나 이 방식은 평가 대상 모델과 채점 모델 사이의 순환적 편향을 내재하며, 반복 실행마다 점수가 변동한다는 한계가 있었다. 이를 보완하기 위해 핵심 탐지 지표의 채점 루프에서 LLM을 배제한 결정론적 오라클 채점기 (Deterministic Oracle Scorer, \texttt{src/agents/oracle\_scorer.py})를 도입하였다. 오라클은 각 데이터셋의 정답 라벨인 공격 범주와 실제 공격 개시 시각을 보유하며, 에이전트의 탐지 결과를 해당 라벨과 비교하여 참 양성 (TP)과 거짓 음성 (FN)을 판정한다.
시스템 평가 초기 단계에서는 각 에피소드별로 에이전트 (LLM) 자신이 산출한 진단 결과를 스스로 $0\sim100$점 척도로 채점하는 \textit{LLM-as-a-judge} 방식을 운용한다. 그러나 이 방식은 가해 모델과 평가 채점 모델 간의 순환적 확증 편향이 개입할 여지가 높고 반복 실행 시 점수가 변동된다는 본질적 한계를 지닌다. 
%
이를 기술적으로 보완하기 위해 핵심 탐지 성능 채점 루프에서 언어 모델을 완전히 배제한 결정론적 오라클 채점기 (deterministic oracle scorer, \texttt{src/agents/oracle\_scorer.py})를 도입한다. 오라클 시스템은 평가 데이터셋에 기록된 정답 라벨인 취약점 범주와 실제 공격 트리거 시각을 독점적으로 보유하며, 에이전트 측의 최종 진단값과 매칭해 참 양성 (TP) 및 거짓 음성 (FN)을 판정한다. 
%
공격 개시 시각과 위협 범주는 공격 주입기가 기록한 결정론적 로그를
ground truth로 사용한다. 반면 임무 훼손 여부와 공격 성공 여부는 위치 이탈, 모드 전환,
웨이포인트 미도달 등 사전에 정의한 임계 기준에 따라 별도로 판정한다.
따라서 공격 개시 시각은 직접 기록된 사건 라벨인 반면, 임무 훼손 여부는 연구자가 사전에 정의한
평가 기준에 의존하는 파생 지표로 구분하여 해석한다.
%%%%%
TP 판정 기준은 \emph{``공격 발생을 정밀하게 인지하고, 방어 에이전트가 보고한 개시 시각이 실제 개시 시점 또는 그 전조 현상 발생 시점으로부터 $2.0\,\text{s}$ 이내에 포착될 것''}으로 정의하였다. 하위 사건 (예: \texttt{mike} 시나리오의 오버라이드 해제 구간) 을 실제 최초 개시 사건으로 오인 탐지할 경우 FN으로 처리한다.

\begin{table}[t!]
\centering
\caption{오라클 정답 라벨 (Ground Truth)}
\label{tab:oracle_gt}
\renewcommand{\arraystretch}{1.2}
\small
\begin{tabularx}{\textwidth}{l l l l X}
\toprule
\textbf{캡처 파일} & \textbf{위협 범주} & \textbf{플랫폼} & \textbf{실제 개시 (s)} & \textbf{비고} \\
\midrule
\texttt{kilo} & \texttt{rc\_override} & UGV & 9.062 & 조향 단일 채널; 부차 사건 \texttt{terminal\_disarm} 25.14\,\text{s} \\
\texttt{lima} & \texttt{datalink\_dos} & Satlink & 18.059 & 전조 현상 \texttt{heartbeat\_stop} 15.0\,\text{s} \\
\texttt{mike} & \texttt{rc\_override} & UAV & 12.048 & 피치 단일 채널; 부차 사건 \texttt{pitch\_pwm\_zero} 34.09\,\text{s} \\
\texttt{ugv\_auto\_cmd} & \texttt{command\_injection} & UGV & 14.160 & AUTO 모드 이탈 강제 (위조 \texttt{DO\_SET\_MODE}) \\
\bottomrule
\end{tabularx}
\end{table}

\subsection{결정론적·재현 가능한 탐지 지표: 소스별 TP/FP/FN/MTTD 정량 분석}
\label{subsec:objective_metrics}

\begin{table}[t!]
\centering
\caption{소스별 결정론적·재현 가능한 탐지 지표 (근거 데이터: \texttt{docs/evidence/oracle\_eval/oracle\_metrics.json})}
\label{tab:oracle_metrics}
\renewcommand{\arraystretch}{1.25}
\scriptsize
\begin{adjustbox}{width=\textwidth}
\begin{tabular}{l l l l l l l l}
\toprule
\textbf{실험 소스군} & \textbf{에피소드} & \textbf{TP} & \textbf{FN} & \textbf{Recall} & \textbf{범주 정확도} & \textbf{평균 절대 개시 오차 (s)} & \textbf{MTTD (정탐, s)} \\
\midrule
\texttt{blind\_manual}      & 6 & 3 & 3 & \textbf{0.50} & 1.0 & 10.004 (max 22.042) & $-2.04$ ($n=3$) \\
\texttt{agentic\_commander} & 3 & 3 & 0 & \textbf{1.00} & 1.0 & 0.420 (max 1.059)   & $-0.286$ ($n=3$) \\
\texttt{real\_game\_campaign} & 4 & 4 & 0 & \textbf{1.00} & 1.0 & 0.365 (max 1.054)   & $-0.162$ ($n=4$) \\
\bottomrule
\end{tabular}
\end{adjustbox}
\end{table}

도구를 결합한 지휘군 (\texttt{agentic\_commander}) 은 \texttt{kilo} 및 \texttt{mike} 시나리오의 실제 개시 시점을 약 $0.1\sim0.2\,\text{s}$ 이내로 정확히 특정한 반면, 도구를 사용하지 않은 도구 미사용 방어군 (\texttt{blind\_manual}) 은 \texttt{mike} 시나리오에서 $34.09\,\text{s}$ 시점의 오버라이드 해제 구간을 최초 공격 시점으로 오인하여 $22.042\,\text{s}$ 지연된 FN을 발생시켰다.

% \noindent \textbf{평가 지표 해석 시의 아키텍처적 유보 조항 (Caveats):} 본 실험에서 재현율 (Recall) 은 시스템의 탐지 능력을 대변하는 대표 지표로 기능할 수 있으나, 정밀도 (Precision) 지표는 정상 음성 데이터가 충분하지 않아 신뢰성있는 지표로 사용될 수 없다. 오탐 (FP) 여부는 공격 개시 이전의 정상 시간 윈도우에 대해서만 제한적으로 검사되었으며, 전체 채점 대상 데이터셋이 공격 (양성) 캡처로만 구성되어 있어 정상 상태 대 공격 상태 간의 교차 특이도 (Specificity) 가 확립되지 않았기 때문이다. 

평균 탐지 시간 (MTTD) 에 나타난 음수 ($-$) 값은 시스템의 오작동이 아닌, 공격의 '전조 현상'에 기반한 조기 탐지 성공을 의미한다. 일례로 \texttt{lima} 시나리오의 경우, D4 도구가 $15.0\,\text{s}$ 부근에서 발생한 하트비트 블랙아웃 전조를 $17.0\,\text{s}$ 시점에 안정적으로 래치함으로써, 최종 라벨링된 RTT 스파이크 공격 개시 시점 ($18.06\,\text{s}$) 보다 빠르게 위협을 포착하였다. 

한편, 평균 복구 시간 (MTTR) 은 메커니즘상 측정 불가 (\texttt{null}) 로 산출된다. 본 평가가 사전에 녹화된 데이터의 개루프 (Open-loop) 재생 방식을 취하므로, 에이전트가 최적의 차단 및 복구 플레이북을 선택하더라도 그것이 실제 기체 제어로 직접 구동 (Actuation) 되지 않기 때문이다. 

오라클 채점 결과는 ``탐지 성공이 곧 임무의 보전을 의미하지는 않는다''는 사실을 실증적으로 명시한다. 실제로 개루프 재생 환경하의 세 시나리오 모두 기체의 임무 능력이 악화된다.

% \subsection{도구 미사용 방어군 vs 도구 결합형 지휘군: 독립 재현과 원래 주장의 정정}
% \label{subsec:blind_vs_tool_commander}

% 초기 대조 실험에서는 도구를 사용하지 않은 도구 미사용 방어군 (Blind Group) 의 단일 공유 BLUE 커맨더 (BLUE commander) 가 \texttt{mike} 시나리오의 위협을 오인하는 현상이 관찰되었다. 실제 위협은 $12.05\,\text{s}$ 시점에 개시된 피치 오버라이드 공격이었으나, 해당 지휘관은 이를 기체의 정상적인 전진 기동으로 판정하여 배제한 후, $34.09\,\text{s}$ 시점에 나타난 오버라이드 해제 구간만을 잔여 이상 징후로 포착하였다. 본 연구진은 이 초기 결과에 기반하여 ``출처 모호성이 극대화된 \texttt{mike} 시나리오는 반드시 탐지 도구가 확보되어야만 해결 가능하다''라는 핵심 결론을 도출한 바 있다. 
% %

% \textbf{그러나 추후 수행된 독립적 재현 실험을 통해 해당 필요성 가설이 성립하지 않음이 증명되었으며, 본 연구진은 이를 다음과 같이 정직하게 정정한다.}

% 본 시스템 아키텍처가 가질 수 있는 내재적 편향을 배제하고 실험 결과의 객관성을 재확증하기 위해, 본 다중 에이전트 시스템을 공유하지 않으며 실험 조건에 대해 완전 블라인드 처리된 별도의 독립적 심판 에이전트 (실험군별 $n=3$) 를 구축하여 재평가를 실시하였다. 구체적인 재채점 결과는 다음과 같다.

% \begin{table}[t!]
% \centering
% \caption{독립 재평가 결과: 도구 미사용 방어군과 도구 결합형 지휘군 간의 비교 검증 ($n=3$, 100점 만점 기준)}
% \label{tab:independent_eval}
% \renewcommand{\arraystretch}{1.2}
% \small
% \begin{tabularx}{\textwidth}{l c c c c}
% \toprule
% \textbf{평가 시나리오} & \textbf{실제 개시 시점 (s)} & \textbf{도구 미사용 방어군 (Blind)} & \textbf{도구 결합형 지휘군 (Agentic)} & \textbf{성능 향상도 (Lift)} \\
% \midrule
% \texttt{mike} (UAV \texttt{rc\_override})  & 12.048 & $93.0 \pm 0.0$  & $95.0 \pm 0.0$ & $+2.0$  \\
% \texttt{kilo} (UGV \texttt{rc\_override})  & 9.062  & $91.0 \pm 1.41$ & $95.0 \pm 0.0$ & $+4.0$  \\
% \texttt{lima} (\texttt{datalink\_dos})       & 18.059 & $82.0 \pm 1.41$ & $94.0 \pm 0.0$ & $+12.0$ \\
% \bottomrule
% \end{tabularx}
% \end{table}

% \noindent 전체 $18$회의 반복 실험 결과, 기계적 진단 정확도는 전수 정답 ($18/18$) 을 기록하며 이는 평가된 세 시나리오 범위에서 두 실험군 모두 안정적으로 위협 범주를 판정했음을 보여준다. 결정적으로 올바르게 통제된 도구 미사용 방어군 역시 \texttt{mike} 시나리오를 $3/3$으로 올바르게 해결하였으며, 공격 개시 시점 추정치의 절대 오차는 $0.0018\,\text{s}$ 수준으로 전체 매트릭스에서 가장 정밀한 축에 속했다.

% 즉, 초기 도구 미사용 방어군에서 관찰된 FN은 후속 독립 재평가에서 재현되지 않았으며 이를 도구 미사용 접근법의 구조적 실패로 일반화할 근거는 부족하다. 결과적으로 ``출처 모호성이 높은 \texttt{mike} 시나리오를 해결하기 위해 도구 아키텍처가 필수적이다''라는 형태의 필요성 주장은 논리적으로 성립하지 않는다.

% 다만 기저 모델을 최고 사양 (\texttt{Opus}) 으로 양측 실험군 모두 엄격히 제한한 통제 변수인 \texttt{lima} 시나리오에서 가장 명확하고 유의미한 성능 향상 ($+12$) 이 관측되었다. 결론적으로 도구 아키텍처는 탐지 성공의 필수조건으로 입증되지는 못했으나, (a) 정량화된 루브릭 점수의 일관된 향장을 유도하며, (b) 자가 보고 방식의 도구 미사용 방어군이 매 반복마다 개시 시각의 변동성을 보이는 것과 달리, 결정론적 탐지기 기반의 명확한 재현성을 보장한다는 점에서 고유한 기술적 가치를 지닌다. 단, 본 심판 프로세스는 시스템 아키텍처를 내재하지 않은 별도의 블라인드 LLM 에이전트에 의해 수행되었으므로, 외부 인간 평가자의 검증과 비교했을 때 잔여 LLM 채점 편향 (Residual Bias) 이 일부 잔존할 수 있다.

\subsection{모델 벤치마크 평가: 평가된 제한 시나리오에서 관찰된 성능 포화와 비용 차이 분석}
\label{subsec:model_bakeoff}

\begin{table}[t!]
\centering
\caption{언어 모델 등급별 벤치마크 평가 결과 (단일 탐지 도구 및 3개 시나리오 통제 환경)}
\label{tab:bakeoff_standard}
\renewcommand{\arraystretch}{1.2}
\small
\begin{tabularx}{\textwidth}{l c c c X}
\toprule
\textbf{모델} & \textbf{탐지 성공 (TP)} & \textbf{평균 수행 시간 (s)} & \textbf{소요 비용 (USD)} & \textbf{비고} \\
\midrule
claude-haiku-4-5       & 3/3 & 28.5 & \textbf{0.0901} & 전 시나리오 교차 확증 달성, 최고 비용 효율성 \\
claude-sonnet-5 (기준) & 3/3 & 27.8 & 0.2258          & 전 시나리오 교차 확증 달성 \\
claude-opus-4-8        & 3/3 & 88.3 & 1.1727          & 전 시나리오 교차 확증 달성 \\
\bottomrule
\end{tabularx}
\end{table}

\noindent 실험 결과, 비교 대상인 세 가지 모델 등급 모두 세 가지 데이터셋 위에서 위협을 정확히 진단하였으며, 복수의 탐지기 신호가 결합하는 교차 확증 상태에 도달하였다. 이는 기 구축된 세 가지 시나리오 환경하에서 모델 규모에 따른 탐지 능력 변별력이 이미 포화 상태에 이르렀음을 의미한다.

따라서 실무적인 관점에서 본 시스템의 유의미한 차별화 요소는 실행 비용으로 귀결된다. 경량 모델인 Haiku는 최고 사양인 Opus 대비 약 $13$배, 중간 등급이자 기준 모델인 Sonnet 대비 약 $2.5$배 수준으로 비용을 절감할 수 있으며, 연산 지연 측면에서는 Sonnet과 대등한 성능을 보여준다.

\subsection{다중 도메인 통합 캠페인 평가: 3개 도메인 조율 Counter-C2 분석}
\label{subsec:real_game_campaign}

\noindent 
% 본 실험에서는 단일하게 조율된 다단계 자율 공격 캠페인을 3개 물리 도메인에 걸쳐 동시 전개하였다. 각 공격 단계는 배포된 도메인별 공격 에이전트 (RED) 가 자체 탑재된 공격 시나리오 라이브러리를 기반으로 언어 모델을 통해 자율적으로 계획하며, 단일 공유 BLUE 커맨더 (\texttt{blue\_commander}) 아키텍처가 이에 대응하는 구도로 실행된다. 전체 캠페인의 상세 시퀀스는 다음과 같다.
본 실험에서는 3개 물리 도메인에 걸쳐 동시 전개되는 다단계 자율 공격 캠페인을 시뮬레이션하였다. 도메인별로 배포된 공격 에이전트(RED)가 자체 탑재된 시나리오 라이브러리와 언어 모델을 활용해 자율적으로 공격을 계획하며, 단일 공유 방어 에이전트(\texttt{blue\_commander})가 이에 대응하는 구조로 진행된다. 전체 캠페인의 상세 시퀀스는 다음과 같다.
\begin{itemize}
    \item \textbf{Stage A (위성/BLOS 도메인):} 데이터링크 서비스 거부 공격을 유도하여 기체의 C2 통신 대역을 압박한다.
    \item \textbf{Stage B (UGV AUTO 킬체인, 총 2단계):} 
    \begin{itemize}
        \item \textbf{Stage B1 (선행 단계):} 위조된 \texttt{DO\_SET\_MODE} (HOLD) 명령을 주입하여 자율 주행 모드로부터 기체를 강제 이탈시킨다.
        \item \textbf{Stage B2 (본 공격 단계):} 모드 이탈로 취약해진 UGV의 조향 제어권을 최종 탈취한다. 제어권 오버라이드 기동 전에 제어 모드를 먼저 변경하도록 설계된 시퀀스는 \textbf{모델이 자율적으로 구성}하였다.
    \end{itemize}
    \item \textbf{Stage C (UAV 도메인):} 피치 축 오버라이드 공격을 가한다.
\end{itemize}

% 전체 캠페인의 실행 과정에서 총 $7$건의 실제 언어 모델 호출 (공격 계획 $4$회, 방어 결정 $3$회) 이 수행되었으며, 소요된 총비용은 $0.6757\,\text{USD}$였다. 기저 모델은 상용 \texttt{claude-sonnet-5} 엔진으로 구동되었으며, 예외 상황으로 인한 규칙 기반 예외 강등 메커니즘은 단 한 차례도 발동되지 않았다.
전체 캠페인 진행 중 총 7회의 언어 모델 호출(공격 계획 4회, 방어 결정 3회)이 발생하였으며, 총 0.6757 USD의 비용이 소요된다. 모든 에이전트는 상용 \texttt{claude-sonnet-5}를 기반으로 구동되며, 모델에 내재된 안전 가드레일로 인한 답변 거부 현상은 관측되지 않았다.

\begin{table}[t!]
\centering
\caption{통합 시뮬레이션 캠페인 기반 단일 공유 BLUE 커맨더 (\texttt{blue\_commander}) 방어 성능 검증 결과}
\label{tab:real_game}
\renewcommand{\arraystretch}{1.2}
\small
\begin{tabularx}{\textwidth}{c l l c c}
\toprule
\textbf{공격 단계 (링크)} & \textbf{대상 시나리오 (캡처)} & \textbf{위협 범주} & \textbf{탐지 여부 (TP)} & \textbf{탐지 소요 시간 (TTD, s)} \\
\midrule
B1 (선행 단계) & \texttt{ugv\_auto\_cmd} & \texttt{command\_injection} & \checkmark & $0.000$ \\
B2             & \texttt{kilo}          & \texttt{rc\_override}       & \checkmark & $0.206$ \\
A              & \texttt{lima}          & \texttt{datalink\_dos}      & \checkmark & $-1.054$ \\
C              & \texttt{mike}          & \texttt{rc\_override}       & \checkmark & $0.202$ \\
\bottomrule
\end{tabularx}
\end{table}

\noindent 
% 실험 결과, 총 \textbf{4/4개의 전체 사이버 킬체인 링크에 대한 방어} 를 달성하였다. 세부적으로 선행 단계인 Stage B1의 경우, D2 레벨의 명령 침입 탐지 시스템이 비정상 시스템 ID (Rogue \texttt{sysid} 200) 및 예기치 않은 모드 전환 명령을 수신 단계에서 즉각적으로 무력화하였다. 이어진 Stage B2의 조향 오버라이드 효과는 D5 도구를 통해 공격 개시 후 $0.206\,\text{s}$ 만에 정밀하게 포착되었다.
실험 결과, 구성된 \textbf{4개의 사이버 킬체인 링크 전체(4/4)에 대한 방어}에 성공하였다. 세부적으로는, 선행 단계인 Stage B1에서 D2 레벨의 명령 침입 탐지 시스템이 비정상 시스템 ID(Rogue \texttt{sysid} 200)와 예기치 않은 모드 전환 명령을 수신 즉시 탐하였다. 이어진 Stage B2의 조향 오버라이드 시도 역시 D5 도구를 통해 공격 개시 0.206초 만에 정밀하게 탐지된다.

% \textbf{기술적 제약 사항:} 본 통합 시뮬레이션 캠페인에서 단일 공유 BLUE 커맨더가 수행한 모든 탐지 활동은 실시간 라이브 폐루프 제어 환경이 아닌, 사전에 독립적으로 기록된 캡처 파일들을 대상으로 수행되었다. 즉, 교차 도메인 간의 유기적인 조율 및 타임라인 시퀀스는 분리 수집된 3개의 시나리오 데이터를 기반으로 논리적으로 재구성한 서사적 아키텍처의 검증이며, 단일하게 동기화되어 실시간 상호작용이 일어나는 다중 기체 라이브 물리 실행 환경을 의미하지 않는다.

\subsection{D5 탐지기의 경험적 탐지 경계 분석}
\label{subsec:defense_robustness_envelope}

\noindent 
% 제어 불변식 탐지기 (D5) 가 가지는 탐지 경계의 정량적 특성을 분석하기 위해, 라이브 SITL 환경에서 오버라이드 제어 입력의 진폭 (PWM 편차 변수, $\mu\text{s}$) 을 매개변수 스윕하며 탐지 소요 시간을 측정하였다. 실험은 기존 수정 전 상태인 기준 (``baseline'') 탐지기와 신규 알고리즘을 부가한 하드닝 (``hardened'') 탐지기 (\texttt{control\_invariant\_hardened.py}) 간의 대조 방식으로 수행되었다.
제어 불변식 탐지기 (D5)의 정량적 탐지 경계 (Detection Envelope)를 분석하기 위해, 라이브 SITL 환경에서 오버라이드 제어 입력 진폭(PWM 편차, μs)을 매개변수로 하여 스윕 (sweep) 테스트를 수행하고 탐지 소요 시간을 측정하였다. 실험은 기존 수정 전 상태인 기준 (``baseline'') 탐지기와 신규 알고리즘을 부가한 하드닝 (``hardened'') 탐지기 (\texttt{control\_invariant\_hardened.py}) 간의 대조 방식으로 진행된다.

\begin{table}[t!]
\centering
\caption{UAV 피치 오버라이드 탐지 경계 분석}
\label{tab:d5_envelope}
\renewcommand{\arraystretch}{1.2}
\small
\begin{tabularx}{\textwidth}{c l c c}
\toprule
\textbf{오버라이드 편차 ($\mu\text{s}$)} & \textbf{실제 발생 물리 효과 (위치 드리프트)} & \textbf{기준 모델 TTD (s)} & \textbf{하드닝 모델 TTD (s)} \\
\midrule
150 & 급격한 위치 드리프트 ($60.87\,\text{m}$) & 0.204 & 0.204 \\
45  & 가시적인 위치 드리프트 ($10.05\,\text{m}$) & 0.204 & 0.204 \\
40  & 완만한 위치 드리프트 ($7.61\,\text{m}$)   & \textbf{10.729} & 0.205 \\
30  & 미세한 위치 드리프트 ($3.61\,\text{m}$)   & \textbf{14.102} & 0.204 \\
20  & 물리적 변위 거의 없음 ($0.01\,\text{m}$)  & \textbf{MISS}   & 2.449 \\
12  & 물리적 변위 없음 ($0.00\,\text{m}$)       & \textbf{MISS}   & 7.561 \\
\bottomrule
\end{tabularx}
\end{table}

\noindent 측정 데이터 분석 결과, 기준 모델은 주입 편차가 $40\,\mu\text{s}$ 이하인 영역 (dev $\le 40\,\mu\text{s}$) 에서 명령 불변식 임계치를 트리거하지 못하는 치명적인 탐지 사각지대를 드러냈다. 이로 인해 기체가 정상 경로를 이미 $3.6\sim7.6\,\text{m}$ 이상 이탈하여 동역학적 물리 드리프트가 가시화된 이후에야 뒤늦게 위협을 포착하거나, $20\,\mu\text{s}$ 이하의 미세 스텔스 주입 공격은 전혀 탐지하지 못하였다(\textbf{MISS}) UGV 도메인의 조향 제어 실험에서도 이와 유사한 탐지 경계 확장 효과 (RC PWM 펄스폭 변조량 기준 하한 임계치 $40\,\mu\text{s} \rightarrow 12\,\mu\text{s}$ 축소) 가 확인된다.
% \noindent \textbf{하드닝 기제에 적용된 두 신호의 기여도 분리 (Ablation Study):} 
추가적으로, 하드닝 기제에 적용된 두가지 신호인 (i) 누설형 단측 CUSUM (Cumulative Sum) 알고리즘 (임계치 미만의 지속적인 서브 임계치 오프셋을 유계 시간 내에 누적 포착) 과 (ii) 기체의 자율 홀드 모드 진입 시 순간 명령 허용 임계치를 기존 $40\,\mu\text{s}$에서 $20\,\mu\text{s}$로 능동 축소하는 적응형 임계치 메커니즘을에 대한 분리 실험 (ablation study)를 진행하였다.
%
소스 코드 분리 실험 결과, CUSUM 알고리즘 단독 적용 시에는 상태 불변식의 지연 반응 공백을 효과적으로 단축시켰으나 적응형 임계치 단독 적용 시에는 데드존 바로 직상위 대역에 대해서만 즉각적인 ($0.2\,\text{s}$) 반응을 보일 뿐 데드존 하위의 미세 신호는 놓치는 한계를 보였다. 즉, 두 신호 메커니즘은 상호 보완적이며 어느 한쪽도 아키텍처적으로 중복 자원이 아님이 규명된다. 

% \noindent \textbf{오탐 (FP) 한계 비용의 객관적 계량.} 
% 총 $14$개의 라이브 정상 상태 비행 윈도우 (UAV 8개, UGV 6개) 세션을 대상으로 검증한 결과, 기존 기준 모델과 하드닝 모델 모두 오탐을 유발하지 않았다. 센서 지터 및 돌풍 환경을 모사한 합성 데이터 $3,000$ 샘플 기반의 교차 검증에서도 피크 CUSUM 누적 수치가 $20.8\sim21.1$ 수준에 머물러, 실제 트립 임계치 ($300$) 대비 약 $14$배 이상의 충분한 구조적 마진을 확보함으로써 오탐 제로를 달성하였다. 그러나 적응형 임계치 메커니즘 도입에 따른 한계 비용은 영 ($0$) 이 아니다. 기체가 \texttt{LOITER} 모드에서 제자리를 유지하기 위해 정상적으로 수행하는 $25\sim30\,\mu\text{s}$ 범위의 미세 제어 조작 (기존 모델이 허용하던 변동 폭) 에 대해서는, 하드닝 시스템하에서 $1\text{ FP}$의 불필요한 오탐 발화가 발생하는 부작용이 관측되었다.

% \noindent \textbf{핵심 구조적 유보 조항:} 본 제어 불변식 탐지기 (D5) 는 시스템의 차단선 역할을 수행할 뿐, 공격 행위 자체를 물리적으로 제어하는 예방기가 아니다. 즉, 하드닝 알고리즘의 도입은 평균 탐지 시간을 획기적으로 낮추고 스텔스성 주입 공백을 폐쇄할 뿐, 후속 완화 제어가 물리적으로 연동되기 전까지는 공격성공률을 독립적으로 낮출 수 없다.

\subsection{실시간 폐루프 능동 대응 메커니즘 검증}
\label{subsec:closed_loop_prevention}

\noindent 앞선 절에서 언급한 ``탐지 성공이 임무 보전과 직결되지 않는다''는 아키텍처적 한계와 ``D5 탐지기는 제어단 액추에이션 전까지 공격성공률을 낮추지 못한다''는 기술적 한계점은, 제어 완화 플레이북이 실제 하드웨어 제어 명령으로 물리적 발송되는 \textbf{별도의 실시간 폐루프 시뮬레이션 실험} 을 통해 해소된다. 

% RC 오버라이드 공격 시나리오 환경하에서 D5 탐지 신호를 실시간 \texttt{BRAKE} 제어 액추에이션과 동적 결합한 결과, 기체의 공격성공률은 기존 3/3 ($1.0$)에서 0/3 ($0.0$)으로 완벽히 통제되었다. 아울러 물리적 피크 변위 (peak deviation) 가 $27.1\,\text{m}$에서 $0.68\,\text{m}$ 수준으로 대폭 축소 ($97.5\%$ 감소) 되었으며, 이때의 평균 탐지 및 제어 소요 시간은 $0.31\,\text{s}$ ($n=3$) 로 수렴하였다. 이를 통해 개루프 재생 환경에서 완전히 파괴되던 기체의 임무 능력이 실시간 폐루프 구조하에서 안정적으로 보전될 수 있음이 증명되었다
RC 오버라이드 공격 상황에서 D5 탐지 신호를 실시간 \texttt{BRAKE} 제어 액추에이션과 동적으로 연동한 결과, 기체의 공격 성공률은 기존 100\%(3/3)에서 0\%(0/3)로 완벽하게 억제된다. 또한, 물리적 최대 변위(Peak Deviation)가 $27.1$m에서 $0.68$m 수준으로 대폭 축소($97.5$\% 감소)되며, 이때 탐지부터 제어까지 소요된 평균 시간은 $0.31$초로 측정된다. 이를 통해 개루프 재생 환경에서는 완전히 상실되던 기체의 임무 능력이, 실시간 폐루프 구조를 통해 안정적으로 보전될 수 있음이 입증된다.

\subsection{연구의 실증적 한계점}
\label{subsec:limitations}

\noindent 본 연구에서 제시하는 모든 실측 지표와 실험적 기여는 다음과 같은 구조적 한계점 내에서 성립하며, 이를 투명하게 밝힌다.

\begin{enumerate}[leftmargin=*]
    \item \textbf{시뮬레이션 환경에 대한 높은 의존성:} 모든 실험은 실제 물리 하드웨어가 아닌 ArduPilot SITL 가상 환경 하에서 수행된다. 데드존 (dead-zone) 폭 등의 수치적 하한은 시뮬레이터 내부의 물리 파라미터에 종속적이다. 그러나 탐지 포락선 (envelope) 의 구조적 특성 (예: 명령 불변식이 특정 $\mu\text{s}$ 임계치 이하에서 탐지 불가 상태가 되거나, 상태 불변식이 데드존 부근에서 지연 및 누락되는 현상) 은 타깃 환경과 무관하게 통계적 유효성을 갖는다.
    
    \item \textbf{소표본 조건 ($n \le 5$):} 본 연구에서 확보한 실험 데이터의 표본 규모는 소규모로 구성된다 (오라클 검증용 소스별 $3\sim6$개 에피소드, 독립 재평가 및 모델 벤치마크 실험군별 각각 $n=3$, 예방 및 강건성 검증 각각 $n=3$). 따라서 본 평가는 방대한 표본에 기반한 엄밀한 통계적 검정력 (statistical power) 의 확보보다는, 실험의 재현성 및 방향성 (부호 안전성) 을 실증적으로 검증하는 데 주안점을 둔다. 
    
    \item \textbf{개루프 제어 환경에 따른 완화 기제의 제약:} 성능 평가 시나리오는 사전에 녹화된 데이터의 개루프 재생 방식을 취하므로, 에이전트가 최적의 차단 및 복구 정책을 도출하더라도 실제 물리적 제어 (actuation) 로 직접 연동되지 않는다. 이로 인해 평균 복구 시간은 측정 불가 (\texttt{null}) 로 산출된다. 완화 기제의 실질적인 차단 능력은 제\ref{subsec:closed_loop_prevention}절의 실시간 폐루프 (live closed loop) 실험을 통해서만 제한적으로 검증된다.
    
    \item \textbf{교차 특이도 및 정밀도 지표의 미확립:} 평가 데이터셋이 대부분 양성(공격) 시나리오로 구성되어 있어, 광범위한 정상 컨텍스트와 대조하는 교차 특이도 및 정밀도 지표를 통계적으로 확립하기 어려웠다. 이에 본 연구는 시스템의 최우선 평가지표로 재현율(Recall)을 채택하였다.
    
    \item \textbf{LLM 기반 평가 체계의 잔여 편향 가능성:} 오라클 채점 루프에서 언어 모델을 완전히 격리하고 외부의 독립적인 블라인드 에이전트를 도입하여 평가의 객관성을 높였으나, 평가의 최종 주체가 인간이 아니며 진단 대상 산출물 역시 LLM 인터페이스를 통해 생성된다는 점에서 기저 파운데이션 모델 공유로 인한 근본적인 잔여 편향이 일부 잔존할 수 있다.
    
    \item \textbf{사이버 킬체인 시나리오의 합성적 특성:} 본 연구에서 수행한 다중 도메인 조정 실험은 독립적인 타임라인에서 개별 수집된 3개의 고유 시나리오 데이터를 바탕으로 복합 공격 킬체인의 성립 가능성을 논리적으로 재구성한 결과이다. 따라서 이는 다수의 기체가 실시간으로 동기화되어 상호작용하는 물리적 라이브 실행 환경을 의미하지 않는다.
\end{enumerate}

\noindent 상기 기술된 한계점들은 도출된 연구 성과와 결론을 약화시키기 위함이 아니며, 본문에서 인용된 모든 실측 수치들이 어떠한 환경적 제약 조건 하에서 과학적 정당성을 확보하는지 그 경계를 명확히 획정하기 위함이다. 본 연구는 실측된 객관적 근거가 확보된 명제만을 주장하며, 메커니즘상 계측 불가능한 영역은 명확히 구분하여 서술한다.